% !TeX encoding = UTF-8
% Español de España. Cargar un único arquetipo desde el archivo principal.
\newcommand{\cvsummary}{Ingeniero mecánico y mecatrónico por NYU Abu Dhabi, con experiencia en I+D de ciencias espaciales en EAU. Combino procesamiento de datos asistido por inteligencia artificial y aprendizaje automático con trabajo práctico en control embebido, diseño de sensores y sistemas robóticos.}

% --- COMPETENCIAS ---
\newcommand{\skills}{
    \section*{COMPETENCIAS TÉCNICAS}
    {\small
    \textbf{Sistemas autónomos y robótica:} Fusión sensorial (IMU), control en tiempo real, adquisición de nubes de puntos 3D y teleoperación, procesamiento de imágenes multiespectrales, MATLAB/Simulink y simulaciones de vuelo.

    \vspace{1pt}
    \textbf{Programación e IA:} Python, C++, MATLAB, JavaScript, PostgreSQL, Git; NumPy, SciPy, conocimientos básicos de PyTorch/fastai, procesamiento de imágenes con aprendizaje automático, fundamentos de visión artificial, procesamiento del lenguaje natural e integración de API, procesos de tratamiento de datos.

    \vspace{1pt}
    \textbf{Sistemas embebidos y electrónica:} Microcontroladores, PWM, ADC/DAC, codificadores en cuadratura, control de motores de corriente continua, máquinas de estados finitos, Arduino, Raspberry Pi, PLC, IoT e integración de hardware y software.

    \vspace{1pt}
    \textbf{Estructuras aeroespaciales y simulación:} CATIA V5, SolidWorks (CSWA), ANSYS (Mechanical, Fluent), COMSOL Multiphysics, análisis por elementos finitos (FEA), dinámica de fluidos computacional (CFD), análisis de vibraciones, diseño ligero, materiales compuestos y fabricación aditiva.

    \vspace{1pt}
    \textbf{Microfabricación y MEMS:} Diseño y fabricación de sensores MEMS piezorresistivos, fotolitografía y diseño de máscaras, deposición de películas delgadas (PVD/pulverización catódica), secuencias de proceso a escala de oblea, selección de materiales, metrología y tolerancias geométricas y dimensionales (GD\&T).

    \vspace{1pt}
    \textbf{Datos espaciales y geoespaciales:} Cartografía planetaria, datos satelitales y multiespectrales, PostGIS y escaneado 3D.

    \vspace{1pt}
    \textbf{Idiomas:} Español nativo, inglés C2, francés C2, italiano A1 y alemán A1.
    }
}

% --- PROYECTOS ---
\newcommand{\vawttitle}{VAWT -- Aprovechamiento de energía eólica urbana}
\newcommand{\vawtdates}{En curso (fecha de inicio no especificada)}
\newcommand{\vawtlocation}{Abu Dabi, EAU}
\newcommand{\vawtdescOne}{Desarrollé un concepto de aprovechamiento de energía eólica urbana mediante el diseño de un aerogenerador de eje vertical orientado al flujo de aire variable generado por el tráfico.}
\newcommand{\vawtdescTwo}{Avancé en la evaluación del diseño aerodinámico mediante simulaciones CFD en ANSYS Fluent, modelando el concepto del aerogenerador antes de tomar decisiones de diseño físico.}

\newcommand{\fsaetitle}{Formula SAE -- Suspensión y chasis de fibra de carbono}
\newcommand{\fsaedates}{Otoño de 2024 -- Actualidad}
\newcommand{\fsaelocation}{Abu Dabi, EAU}
\newcommand{\fsaedescOne}{Desarrollé el diseño y la implementación de la suspensión y del chasis de fibra de carbono de un vehículo, combinando modelado CAD y análisis por elementos finitos. Participación en Formula SAE en la Universidad de Oviedo y NYUAD.}
\newcommand{\fsaedescTwo}{Evalué el comportamiento de la suspensión mediante análisis cinemáticos en SolidWorks y ANSYS para estudiar sus prestaciones dinámicas.}

\newcommand{\missiletitle}{Simulador de dinámica de interceptación de misiles}
\newcommand{\missilelocation}{Besançon, Francia}
\newcommand{\missiledescOne}{Desarrollé en MATLAB una simulación 3D de interceptación de misiles mediante guiado por navegación proporcional, integración numérica y visualización de trayectorias en tiempo real con detección de impacto.}

\newcommand{\tactitle}{Mininevera -- Sistema termoeléctrico TEC de dos etapas}
\newcommand{\tacdates}{Primavera de 2025}
\newcommand{\taclocation}{Abu Dabi, EAU}
\newcommand{\tacdescOne}{Desarrollé un sistema térmico de refrigeración termoeléctrica de dos etapas sin refrigerante, utilizando módulos TEC de estado sólido en lugar de un circuito frigorífico.}
\newcommand{\tacdescTwo}{Integré sensores e instrumentación térmica en tiempo real en el prototipo para supervisar el comportamiento del sistema de refrigeración.}

\newcommand{\orbitaltitle}{OrbitalGrid -- Sistema de información geográfica espacial}
\newcommand{\orbitaldates}{06/2024 - Actualidad}
\newcommand{\orbitallocation}{Abu Dabi, EAU}
\newcommand{\orbitaldescOne}{Desarrollé un sistema de datos geoespaciales mediante un proyecto de cartografía aeroespacial basado en SQL, integrando escaneado 3D y datos satelitales.}
\newcommand{\orbitaldescTwo}{Implementé la gestión de datos espaciales en PostgreSQL, estructurando la información geoespacial para aplicaciones cartográficas.}

\newcommand{\softrobtitle}{Exoesqueleto robótico para rehabilitación tras un ictus}
\newcommand{\softrobdates}{04/2024 - 05/2024}
\newcommand{\softroblocation}{Abu Dabi, EAU}
\newcommand{\softrobdescOne}{Diseñé un mecanismo de robótica blanda para rehabilitación mediante un exoesqueleto de dedos inspirado en el origami y accionado por hilos de aleación con memoria de forma.}
\newcommand{\softrobdescTwo}{Integré materiales flexibles y elementos activos en un conjunto de polímero Dragon Skin e hilos SMA, diseñado para reproducir la función muscular de los dedos.}

\newcommand{\exoplanettitle}{Simulador de tránsitos de exoplanetas}
\newcommand{\exoplanetdates}{Primavera de 2025}
\newcommand{\exoplanetlocation}{Abu Dabi, EAU}
\newcommand{\exoplanetdescOne}{Desarrollé un simulador de tránsitos de exoplanetas en Python basado en relaciones de mecánica orbital y física estelar.}
\newcommand{\exoplanetdescTwo}{Generé datos de observación sintéticos en forma de series temporales de curvas de luz, calculando el comportamiento de los tránsitos a partir de parámetros astrofísicos.}

\newcommand{\windTurbinetitle}{Aerogenerador marino flotante -- Estabilización activa}
\newcommand{\windTurbinedates}{05/2026 - 06/2026}
\newcommand{\windTurbinelocation}{Besançon, Francia}
\newcommand{\windTurbinedescOne}{Logré la estabilización en lazo cerrado de la góndola en ensayos de laboratorio en canal de oleaje mediante control PID de dos ejes en Arduino con realimentación de orientación de un MPU6050.}
\newcommand{\windTurbinedescTwo}{Diseñé mecánicamente una plataforma marina semisumergible de tres columnas con un sistema de control para estabilizar un prototipo de aerogenerador a escala 1:1000, analizando su cabeceo e inclinación mediante un sensor MPU6050.}

\newcommand{\carbonmetertitle}{CARBONMETRE -- Ecodiseño y medición de carbono}
\newcommand{\carbonmeterdates}{09/2025 - 12/2025}
\newcommand{\carbonmeterlocation}{Besançon, Francia}
\newcommand{\carbonmeterdescOne}{Apliqué ingeniería de medida a la evaluación ambiental mediante un proyecto de medición de carbono que combinaba metrología, fabricación y ecodiseño.}
\newcommand{\carbonmeterdescTwo}{Relacioné la fabricación con el análisis de sostenibilidad mediante la aplicación de principios de ingeniería de fabricación industrial al proyecto.}

\newcommand{\hvactitle}{Acondicionamiento térmico de una habitación de residencia -- Refrigeración pasiva y evaporativa}
\newcommand{\hvacdates}{06/2026}
\newcommand{\hvaclocation}{Besançon, Francia}
\newcommand{\hvacdescOne}{Logré una reducción de temperatura declarada de hasta 3 \textdegree{}C mediante una pantalla solar exterior y un prototipo de refrigeración con agua congelada y ventilador en una habitación de 9 m$^2$.}
\newcommand{\hvacdescTwo}{Reduje la exposición solar directa en las horas punta de 14:00 a 21:00 mediante capas exteriores de papel de aluminio reflectante, EPS de 1 cm y cartón ondulado.}

\newcommand{\controlBoxtitle}{Instalación eléctrica integral -- Sistemas para eventos de mediana escala}
\newcommand{\controlBoxdates}{06/2023 - 08/2023}
\newcommand{\controlBoxlocation}{Ciudad de México, México}
\newcommand{\controlBoxdescOne}{Reduje un 30\% el tiempo de montaje del sistema mediante el diseño de un cuadro de control de 12 entradas para instalaciones eléctricas de eventos de mediana escala.}
\newcommand{\controlBoxdescTwo}{Desarrollé un sistema eléctrico coordinado basado en PLC para eventos, integrando funciones de control adaptadas a un entorno de funcionamiento en directo.}

\newcommand{\hyperlooptitle}{NYC Hyperloop -- Vehículo a escala reducida}
\newcommand{\hyperloopdates}{01/2024 - 05/2024}
\newcommand{\hyperlooplocation}{Nueva York, EE. UU.}
\newcommand{\hyperloopdescOne}{Lideré el desarrollo mecánico de un sistema de transporte avanzado como responsable del subequipo mecánico, codesarrollando un Hyperloop a escala reducida para un concepto de trayecto Nueva York--Long Island.}
\newcommand{\hyperloopdescTwo}{Evalué el comportamiento del transporte electromagnético mediante simulaciones de levitación magnética y el uso de COMSOL, ANSYS y SolidWorks para el diseño del sistema.}

\newcommand{\dronetitle}{Cuadricóptero funcional -- Estructura ligera}
\newcommand{\dronedates}{03/2024 - 05/2024}
\newcommand{\dronelocation}{Nueva York, EE. UU.}
\newcommand{\dronedescOne}{Reduje un 40\% el peso del cuadricóptero mediante el diseño de una estructura impresa en 3D en PLA y TPU y su integración en un vehículo funcional.}
\newcommand{\dronedescTwo}{Evalué la resistencia al impacto de la estructura del cuadricóptero mediante ensayos de caída contrastados con los criterios de rendimiento durante el desarrollo.}

\newcommand{\trentxwbtitle}{Rolls-Royce Trent XWB -- Modelo 3D funcional}
\newcommand{\trentxwbdates}{01/2024 - 05/2024}
\newcommand{\trentxwblocation}{Nueva York, EE. UU.}
\newcommand{\trentxwbdescOne}{Codesarrollé un modelo detallado del sistema de propulsión Trent XWB, con una representación 3D funcional completa en SolidWorks y Fusion 360 y evaluación estructural y de flujo en ANSYS.}
\newcommand{\trentxwbdescTwo}{Evalué el comportamiento aerodinámico mediante simulaciones y herramientas de análisis de flujo en ANSYS.}

\newcommand{\liostitle}{LIOS -- Asistente de voz con IA}
\newcommand{\liosdates}{12/2024 - Actualidad}
\newcommand{\lioslocation}{Abu Dabi, EAU}
\newcommand{\liosdescOne}{Desarrollé un sistema interactivo con IA mediante un asistente de voz en Python que combina procesamiento del lenguaje natural, datos en tiempo real e integración de API.}
\newcommand{\liosdescTwo}{Habilité interacciones sensibles al contexto mediante memoria contextual, manteniendo información relevante entre las tareas del asistente.}

\newcommand{\fitpaltitle}{FitPal -- Dispositivo inteligente para entrenamiento físico}
\newcommand{\fitpaldates}{09/2021 - 12/2021}
\newcommand{\fitpallocation}{Abu Dabi, EAU}
\newcommand{\fitpaldescOne}{Codesarrollé un producto de sensorización embebida basado en M5Stack: un dispositivo inteligente de entrenamiento físico para analizar ejercicios en tiempo real.}
\newcommand{\fitpaldescTwo}{Conecté el dispositivo a una interfaz móvil mediante funciones de transferencia de datos IoT e integración de Arduino IoT App.}

\newcommand{\covidtitle}{Cabina de desinfección automatizada}
\newcommand{\coviddates}{09/2020 - 02/2021}
\newcommand{\covidlocation}{Ciudad de México, México}
\newcommand{\coviddescOne}{Desarrollé un prototipo automatizado de cabina de desinfección basado en Arduino, integrando electrónica de control y un sistema de nebulización.}
\newcommand{\coviddescTwo}{Implementé comunicación remota mediante enlaces inalámbricos en el sistema de control del prototipo.}

\newcommand{\rovertitle}{Rover teleoperado -- Escaneado espacial 3D}
\newcommand{\roverdates}{09/2026 - Actualidad}
\newcommand{\roverlocation}{Oviedo, España}
\newcommand{\roverdescOne}{Desarrollo una plataforma móvil de sensores que combina control embebido, adquisición de nubes de puntos 3D y visualización remota, con el escaneado espacial teleoperado de principio a fin como objetivo de integración.}
\newcommand{\roverdescTwo}{Desarrollo la teleoperación del rover mediante la integración del control embebido en el sistema móvil, con el funcionamiento por órdenes remotas como hito objetivo.}

\newcommand{\microsensortitle}{Microsensor piezorresistivo de presión de detonación}
\newcommand{\microsensordates}{05/2026 - 06/2026}
\newcommand{\microsensorlocation}{Besançon, Francia}
\newcommand{\microsensordescOne}{Reduje las iteraciones de dimensionamiento del sensor mediante la automatización de cálculos geométricos y de máscaras en una herramienta Python que impone directamente restricciones de K, $\Delta R/R_0$ y resistividad.}
\newcommand{\microsensordescTwo}{Lideré la selección de materiales bajo restricciones para un diseño de sensor a escala de oblea con interfaz de 50 $\Omega$ y resistencia a presiones de 10 GPa, equilibrando requisitos eléctricos, mecánicos y económicos.}

% --- EXPERIENCIA ---
\newcommand{\carbonorg}{Energy Innovation Capital}
\newcommand{\carbontitle}{Analista de mercado de captura de carbono}
\newcommand{\carbondates}{10/2023 - 01/2024}
\newcommand{\carbonlocation}{Abu Dabi, EAU}
\newcommand{\carbonbullets}{%
  \item Apoyé decisiones de inversión mediante recomendaciones sobre empresas emergentes de captura de carbono, analizando tecnologías CCUS para un fondo de capital riesgo de Abu Dabi.
  \item Transformé avances científicos y técnicos en captura de carbono en investigación útil para la toma de decisiones y presentaciones a analistas sénior.
  \item Analicé la incertidumbre comercial y los riesgos de inversión mediante el estudio de tendencias de mercado y factores tecnológicos del sector CCUS.
  \item Relacioné la investigación en ingeniería con la estrategia de inversión mediante recomendaciones que combinaban resultados técnicos y análisis de datos de mercado.
}

\newcommand{\spacelaborg}{NYU Abu Dhabi, Centro de Ciencias Espaciales}
\newcommand{\spacelabtitle}{Ingeniero e investigador de I+D -- Plataforma de ciencias planetarias}
\newcommand{\spacelabdates}{05/2025 - 09/2025}
\newcommand{\spacelablocation}{Abu Dabi, EAU}
\newcommand{\spacelabbullets}{%
  \item Mejoré un 40\% la eficiencia del proceso de tratamiento de datos mediante la optimización de compresión y reproyección, con corrección de imágenes asistida por aprendizaje automático sobre conjuntos de datos a escala planetaria.
  \item Desplegué una plataforma de cartografía científica mediante una \href{https://romvio.github.io/mars-viewer/}{aplicación web y móvil}, integrando imágenes multiespectrales de Emirates Mars Mission HOPE con CSS/JavaScript; posteriormente la presenté a la App Store de Apple en EAU, ampliando el atlas planetario más allá de su versión web.
  \item Integré observaciones de múltiples imágenes en mapas planetarios mediante calibración de solapamientos e implementación de proyecciones cilíndrica y Robinson.
}

\newcommand{\geoorg}{Global Education NYUAD}
\newcommand{\geotitle}{Coordinador de Global Education -- Estudiantes visitantes}
\newcommand{\geodates}{12/2025 - 05/2025}
\newcommand{\geolocation}{Abu Dabi, EAU}
\newcommand{\geobullets}{%
  \item Organicé actividades interculturales en EAU, incluidas salidas académicas y culturales para estudiantes visitantes, mediante la coordinación de actividades internacionales de NYU Abu Dhabi.
  \item Coordiné las operaciones del equipo mediante el liderazgo de los coordinadores y la planificación logística para grupos diversos de estudiantes.
  \item Contribuí a la seguridad de los grupos mediante la planificación de protocolos y la incorporación de precauciones operativas en las salidas.
  \item Mantuve la comunicación entre las partes implicadas como enlace para estudiantes visitantes, conectándolos con los servicios de apoyo universitario.
}

\newcommand{\geoOneorg}{Global Education NYUAD}
\newcommand{\geoOnetitle}{Coordinador de Global Education -- Estudiantes de primer curso}
\newcommand{\geoOnedates}{09/2026 - 12/2024}
\newcommand{\geoOnelocation}{Abu Dabi, EAU}
\newcommand{\geoOnebullets}{%
  \item Facilité la incorporación de estudiantes de primer curso mediante sesiones de orientación e introducciones estructuradas a la vida universitaria.
  \item Ofrecí orientación continuada entre estudiantes mediante mentoría y asesoramiento durante su transición a la universidad.
  \item Comuniqué información sobre instalaciones y recursos universitarios mediante visitas al campus para nuevos estudiantes.
  \item Promoví la integración en la comunidad mediante la planificación de eventos académicos y culturales de inclusión para estudiantes de primer curso.
}

\newcommand{\hcorg}{Centro de Salud de NYUAD}
\newcommand{\hctitle}{Empleo de estudiante -- Responsable de redes sociales}
\newcommand{\hcdates}{09/2024 - 05/2025}
\newcommand{\hclocation}{Abu Dabi, EAU}
\newcommand{\hcbullets}{%
  \item Promoví el uso de recursos de bienestar mediante campañas de salud dirigidas al alumnado y adaptadas a la comunidad de NYU Abu Dhabi.
  \item Elaboré materiales de comunicación digital y contenidos para redes sociales destinados a promover iniciativas de salud.
  \item Actué como punto de contacto para estudiantes, asumiendo la recepción de consultas y coordinando respuestas del centro de salud.
  \item Contribuí a la resolución de incidencias mediante la gestión de consultas, combinando comunicación y apoyo directo al alumnado.
}

\newcommand{\alpesorg}{Alpes Technologies México}
\newcommand{\alpestitle}{Ingeniero de proyectos júnior -- SGA y sistemas eléctricos}
\newcommand{\alpesdates}{06/2024 - 09/2024}
\newcommand{\alpeslocation}{Ciudad de México, México}
\newcommand{\alpesbullets}{%
  \item Aumenté un 60\% la eficiencia operativa y de fabricación mediante la sustitución de la gestión manual de existencias por un sistema propio de gestión de almacenes en HTML/JavaScript y Google Drive.
  \item Implementé trazabilidad industrial mediante seguimiento de números de lote y actualizaciones bidireccionales de existencias en tiempo real, integrando lectura QR, acceso por roles e informes automatizados.
  \item Mejoré el rendimiento eléctrico industrial, con factores de potencia de hasta el 99\%, reducciones del 15\% en costes energéticos y del 25\% en paradas, mediante análisis de subestaciones e instalación de baterías de condensadores.
  \item Realicé el diagnóstico eléctrico de un complejo hotelero tras un incendio mediante mediciones in situ, elaborando esquemas unifilares y recomendaciones correctivas para la alta dirección.
}

\newcommand{\epigsaorg}{EPIGSA}
\newcommand{\epigsatitle}{Asistente de ingeniería de datos -- SIG e hidráulica}
\newcommand{\epigsadates}{05/2024 - 07/2024}
\newcommand{\epigsalocation}{Ciudad de México, México}
\newcommand{\epigsabullets}{%
  \item Reduje un 25\% los errores de datos mediante la automatización de controles de calidad para información de infraestructuras municipales de tratamiento de aguas.
  \item Integré información geoespacial de infraestructuras mediante la consolidación y validación de datos SIG en PostgreSQL para proyectos de tratamiento de aguas.
  \item Mejoré la fiabilidad de los datos de ingeniería mediante validación SIG y comprobaciones de información de infraestructuras durante su integración en bases de datos.
  \item Relacioné la ingeniería de bases de datos con infraestructuras físicas mediante integración espacial y control de calidad automatizado para proyectos de tratamiento de aguas en PostgreSQL.
}

\newcommand{\solarshiporg}{Solar Ship}
\newcommand{\solarshiptitle}{Ingeniero en prácticas de estrategia y diseño de movilidad aérea}
\newcommand{\solarshipdates}{05/2023 - 08/2023}
\newcommand{\solarshiplocation}{Toronto, Canadá}
\newcommand{\solarshipbullets}{%
   %\item Desarrollé un concepto de infraestructura de logística aérea para zonas remotas de Zambia y Tanzania, atendiendo las necesidades de acceso de comunidades sin conexión a la red.
  \item Contribuí al desarrollo de sistemas aeroespaciales eléctricos mediante el diseño de un sistema aeronáutico totalmente eléctrico para la iniciativa Peace+Freedom, orientado a la conectividad de comunidades remotas de Zambia y Tanzania sin conexión a la red.
  \item Relacioné el diseño del vehículo con su contexto operativo en una aplicación de logística remota, considerando la infraestructura junto con el sistema aéreo.
  \item Integré requisitos técnicos y sociales en un proyecto multinacional de movilidad sostenible, combinando diseño aeroespacial y objetivos de acceso de las comunidades.
}

\newcommand{\crousorg}{CROUS Besançon Bourgogne Franche-Comté}
\newcommand{\croustitle}{Empleo de estudiante -- Restauración universitaria}
\newcommand{\crousdates}{10/2024 - Actualidad}
\newcommand{\crouslocation}{Besançon, Francia}
\newcommand{\crousbullets}{%
  \item Mantuve la continuidad del trabajo de cocina en horas punta mediante apoyo al servicio universitario de restauración de gran volumen.
  \item Apliqué procedimientos de limpieza e higiene de cocina, contribuyendo a la disciplina de proceso y a la seguridad alimentaria.
  \item Contribuí a la continuidad operativa mediante reposición de suministros y asistencia durante el servicio de comidas.
  \item Colaboré en tareas coordinadas de cocina y servicio, adaptándome a necesidades operativas con plazos ajustados.
}

\newcommand{\rallyorg}{Campeonato Mundial de Rally -- Guanajuato}
\newcommand{\rallytitle}{Comisario técnico oficial de la FIA}
\newcommand{\rallydates}{11/2019 - 04/2020}
\newcommand{\rallylocation}{Guanajuato, México}
\newcommand{\rallybullets}{%
  \item Realicé inspecciones técnicas reglamentarias y el marcaje oficial de neumáticos en el Campeonato Mundial de Rally de la FIA, aplicando los procedimientos de inspección de la prueba.
  \item Contribuí al cumplimiento reglamentario de los vehículos de competición mediante tareas de verificación técnica.
  \item Trabajé con plazos críticos junto a comisarios responsables y pilotos, coordinando inspecciones en un entorno de competición de alta presión.
  \item Apliqué procedimientos establecidos de marcaje y verificación en trabajos de inspección descritos conforme a normas ISO 9001.
}

\newcommand{\papaloteorg}{Papalote Museo del Niño}
\newcommand{\papalotetitle}{Auxiliar de atención al público}
\newcommand{\papalotedates}{05/2019 - 07/2019}
\newcommand{\papalotelocation}{Ciudad de México, México}
\newcommand{\papalotebullets}{%
  \item Traduje conceptos técnicos para personas no especializadas mediante visitas a exposiciones de ciencia y tecnología, adaptando las explicaciones a niños y adultos.
  \item Contribuí a mejoras implantadas en la experiencia de los visitantes mediante propuestas prácticas durante el trabajo en el museo.
  \item Atendí a visitantes de diferentes edades, ajustando la comunicación a sus necesidades.
  \item Relacioné contenidos técnicos con la participación del público mediante orientación práctica en las exposiciones y explicaciones accesibles de ciencia.
}